% Replacement Chapter 3 for the supplied thesis project.
% Implementation claims were checked against artyomzifir/ViKi, main commit
% 0fdbc02f8f68d82322d32fa9151097034593d30c (2026-10-07).
% This is a source inspection, not a rerun of hardware or benchmark experiments.
% Existing thesis labels are retained where their corrected meaning permits.
% No new bibliography keys are required beyond the revised references.bib.

\chapter{Design and Methodology}
\label{ch:method}

\section{System Overview and Implementation Boundary}
\label{sec:overview}

\viki\ is designed as an offline pipeline that separates human-demonstration
capture, motion reconstruction, robot-specific retargeting, and synthetic
observation generation. Recording a human episode does not require the
demonstrator to control a robot. Subsequent processing may use the complete
episode, rather than meeting the latency requirements of teleoperation.
This permits temporal fitting and trajectory-level optimisation, while making
their processing cost and use of future observations explicit.

The target workflow has two branches after capture. The motion branch
estimates hand and object trajectories and derives robot joint and gripper
commands. The appearance branch reconstructs reusable scene assets and
renders the resulting robot motion from selected cameras. Their outputs
must agree on coordinate frames, timestamps, object identities, and robot
state. Figure~\ref{fig:pipeline} shows this design.

\begin{figure}[htbp]
\centering
\begin{tikzpicture}[
  node distance=7mm and 10mm,
  box/.style={draw,rounded corners=1pt,align=center,font=\footnotesize,
              minimum height=11mm,text width=40mm,inner sep=3pt},
  proposed/.style={box,dashed},
  arr/.style={-{Latex[length=2mm]},thick}]
\node[box] (cap) {Calibrated RGB-D capture};
\node[box,below left=of cap] (motion) {Hand and rigid-object motion};
\node[proposed,below right=of cap] (scene) {3DGS scene assets};
\node[box,below=of motion] (ik) {Embodiment-specific retargeting};
\node[proposed,below=of scene] (render) {Robot and scene rendering};
\node[proposed,below right=of ik] (data) {Checked observation--action episodes};
\draw[arr] (cap) -- (motion);
\draw[arr] (cap) -- (scene);
\draw[arr] (motion) -- (ik);
\draw[arr] (motion) -- (scene);
\draw[arr] (scene) -- (render);
\draw[arr] (ik) -- (render);
\draw[arr] (ik) -- (data);
\draw[arr] (render) -- (data);
\end{tikzpicture}
\caption{Target dataset-generation workflow. Solid boxes have corresponding
processing code in the inspected revision; dashed boxes denote the proposed
3DGS and synthetic-dataset extension. An implemented stage does not imply
that its accuracy or downstream task performance has been established.}
\label{fig:pipeline}
\end{figure}

The implementation description is tied to repository revision
\texttt{0fdbc02} of 7 October 2026.\footnote{Inspected source:
\href{https://github.com/artyomzifir/ViKi/tree/0fdbc02f8f68d82322d32fa9151097034593d30c}{ViKi repository at the pinned revision}.}
At that revision, capture, hand processing, rigid-object modelling,
trajectory-level inverse kinematics (IK), object-relative diagnostics, and
trajectory export have code implementations. An experimental cube-specific
grasp planner also exists. The 3DGS reconstruction and robot-compositing
stages described below are proposed extensions; the existing point-cloud
viewer and image utilities do not implement them. Physical robot replay
remains a stub, and the optional LeRobot writer is partial.

Stage outputs are stored separately so that a failed downstream stage does
not require recapturing the demonstration. The principal current artifacts
are raw RGB-D recordings, per-camera landmarks in \texttt{rec.npz}, prepared
hand trajectories in \texttt{cln.npz}, rigid-object models and tracks in
\texttt{object\_models.npz}, and robot plans in \texttt{plan.h5}.
Object-relative tracks are an additional artifact computed after IK.
The proposed renderer consumes these estimates; it does not turn them into
independently measured robot demonstrations.

\section{Capture Geometry and Coordinate Frames}
\label{sec:hardware}

\subsection{Hardware and workspace}

The local reference configuration uses two Azure Kinect RGB-D cameras
viewing a bounded tabletop workspace. Camera placement should provide
overlapping coverage of the hand and manipulated objects, while avoiding
nearly parallel viewing rays in the interaction region. More views can
reduce particular occlusions, but their benefit depends on placement,
calibration, detector errors, and depth quality. Neither the number of
cameras nor nominal sensor specifications establish reconstruction accuracy
\cite{mskinect,tolgyessy,bulee}.

The acquisition layer also contains a RealSense backend. Supporting a
backend is distinct from validating a particular mixed-camera rig. Every
experiment must record the devices, colour and depth modes, exposure
settings, frame rate, synchronization mode, and calibration used. Camera
motion invalidates the corresponding fixed-rig calibration. The robot is
introduced as a kinematic model during offline processing; a physical robot
is not required for human capture.

\subsection{Calibration and frame convention}
\label{sec:calib}

Let $\bT^A_B$ map coordinates from frame $B$ to frame $A$.
The implementation distinguishes a camera-rig frame $G$, a ChArUco
calibration/workspace frame $C$, and the base frame $B_m$ of target robot
$m$. The rig frame is tied to a reference camera. Inter-camera extrinsics
and the workspace anchor are separate calibration quantities, rather than
two names for the same world frame.

For a point observed by camera $k$, the complete transformation is
\begin{equation}
  \begin{bmatrix}\bx^C\\1\end{bmatrix}
  = \bT^C_G\,\bT^G_k
    \begin{bmatrix}\bx^k\\1\end{bmatrix}.
  \label{eq:extrinsic}
\end{equation}
The camera geometry includes intrinsics and distortion or factory ray
models, together with the relationship between colour and depth sensors.
ChArUco observations provide constraints for the rig solve and workspace
anchor \cite{opencvcharuco}. In the Azure Kinect path, factory-calibrated
rays are used to place colour observations in an ideal pinhole plane for
the multi-pose solve. This avoids treating raw distorted pixels as if they
already followed that ideal model.

The saved calibration must retain the board dimensions and marker
dictionary actually used. A numerically successful calibration solve or
an automatic validation verdict is insufficient evidence of metric
accuracy throughout the task volume. Reprojection residuals, cross-view
surface agreement, and an independent dimension or distance check answer
different questions and should be reported separately.

For retargeting, the configured robot-base pose is $\bT^C_{B_m}$, giving
$\bT^{B_m}_H=(\bT^C_{B_m})^{-1}\bT^C_H$ for a human target frame $H$.
The current code accepts a full base translation and orientation from
configuration. It does not estimate this pose through a physical hand--eye
calibration procedure. The transform therefore defines a virtual placement
for geometric evaluation. Use on hardware would require an independently
verified registration of the actual robot base. No automatic workspace
rescaling or recentering is part of this transformation.

\subsection{Synchronisation}
\label{sec:sync}

The recorder stores both device timestamps and host timestamps. Host time
is derived from a monotonic clock,
\begin{equation}
  \tau_{\mathrm{host},\mu\mathrm{s}}
    = \left\lfloor\frac{t_{\mathrm{monotonic},\mathrm{ns}}}{1000}\right\rfloor,
  \label{eq:hostclock}
\end{equation}
and is used to select nearby buffered frames around a recording tick.
The host-time acceptance window is a stale-frame guard; it does not certify
simultaneous exposure.

For a configured hardware-synchronised Kinect rig, the implementation
checks master/subordinate roles and timestamp-offset stability before
recording. It then matches subordinate device timestamps to the master,
allowing for the measured offset. The inspected revision uses a
$500\,\mu\mathrm{s}$ residual tolerance for this device-timestamp check.
That configured tolerance is not a measured bound on optical exposure
alignment or on the reconstruction error of an episode. Other acquisition
configurations use host-based grouping unless a verified device-time
relationship is available.

Offline multiview reconstruction groups observations by recorded
synchronised frame index. Different cameras' host timestamps need not be
identical within that group. An optional image-landmark time-alignment
step is a distinct processing choice and must be recorded in the profile.
Local synchronization assessment should report device-time residuals,
host-arrival offsets, dropped groups, and drift over the recording. An
external benchmark cannot validate the timing of the local camera rig.

\section{Hand Reconstruction and Temporal Preparation}
\label{sec:fusion}

\subsection{Image observations and depth evidence}

The working MediaPipe backend supplies 21 hand landmarks per detected hand
in each colour image. Each observation contains image coordinates, a
detector score, and, where available, a local depth measurement and depth
spread. Detector confidence is evidence used by the estimator, not a
calibrated probability of 3D correctness. Depth lifting uses the camera's
colour--depth geometry and converts measurements to metres.

Per-camera observations and depth-lifted trajectories are retained before
cross-camera fusion. This permits comparison of triangulation with the
legacy depth-based path without rerunning the detector. Missing detections
and invalid depth are recorded explicitly rather than replaced by zero
coordinates.

\subsection{Calibrated multiview triangulation}

The named multiview profiles use triangulation as their primary fusion
method. For each landmark and frame, camera pairs provide direct linear
triangulation hypotheses. A hypothesis must lie in front of both cameras
and satisfy a minimum angle between its viewing rays. Candidate hypotheses
are ranked by reprojection-inlier count, then inlier detector support and
reprojection error. The best supported hypothesis is refined over its
inlier views.

For the angular-residual mode, let $s_k$ be the detector score, $f_k$ the
camera focal scale, $d_k$ a valid depth observation, and $z_k(\bx)$ the
predicted camera-axis depth. The residual for view $k$ contains
\begin{equation}
\begin{aligned}
  \mathbf r_k(\bx)&=
    \begin{bmatrix}
      \sqrt{s_k}\,[\pi_k(\bx)-\bu_k]/f_k\\[1mm]
      \sqrt{\lambda_d w^d_k}\,
        [z_k(\bx)-d_k-\delta_i]/\max(z_k(\bx),\varepsilon)
    \end{bmatrix},\\[2mm]
  \hat\bx_{i,t}&=\argmin_{\bx}
      \sum_{k\in\mathcal K^{\mathrm{in}}_{i,t}}\sum_a
      \rho\!\left([\mathbf r_k(\bx)]_a\right).
\end{aligned}
  \label{eq:triangulate}
\end{equation}
Here $\pi_k$ projects into the undistorted image convention used by the
solver. The depth row is omitted when depth is invalid. Its weight decreases
with local depth spread; $\delta_i$ is a configurable surface-to-joint
depth offset. The robust loss acts on scalar residual components. The
legacy pixel-residual mode has a different scaling and is preserved as a
separate profile choice.

Triangulation does not remove range-dependent uncertainty. Small ray
intersection angles, inaccurate calibration, motion between views, and
consistent image-localisation errors can all produce an inaccurate 3D
point. Its geometric quality score is
\begin{equation}
  q_{i,t}=\frac{n_{\mathrm{in}}}{n_{\mathrm{usable}}}
    \left[1-\frac{\bar e}{2\epsilon_{\mathrm{in}}}\right]_{0}^{1}
    \left[\frac{\theta_{\mathrm{best}}}{\theta_{\mathrm{ref}}}\right]_{0}^{1},
  \label{eq:quality}
\end{equation}
where $[\cdot]_0^1$ denotes clipping, $\bar e$ is mean inlier reprojection
error, and $\theta_{\mathrm{best}}$ is the largest inlier-pair ray angle.
Profiles may additionally multiply this quantity by mean inlier detector
score. Quality depends on the chosen gates and is not an independent
accuracy measurement. With fewer than two usable views, this triangulation
path leaves a gap.

\subsection{Legacy fusion, interpolation, and articulated fitting}

The alternative \texttt{xyz\_mean} path combines finite depth-lifted points
on a common time grid:
\begin{equation}
  \hat\bx_{i,t}
    =\frac{\sum_k w^k_{i,t}\bx^k_{i,t}}{\sum_k w^k_{i,t}}.
  \label{eq:fusion}
\end{equation}
Its current weights come from the supplied per-camera confidence arrays.
The implementation does not compute the previously proposed product of
detector confidence, depth-sensor confidence, inverse squared range, and
surface-incidence angle. This legacy mode can operate on single-view depth
evidence, but it is not a per-landmark automatic fallback inside the strict
named triangulation profiles. A custom run that falls back because a
triangulation artifact is missing must be identified as such.

After fusion, missing coordinates may be interpolated and the landmark
sequence smoothed with a Savitzky--Golay filter. These are operations on
landmark coordinates, not cubic interpolation on $\SE$. The inspected
default profile, \texttt{stable-fused-hand-v4}, uses linear fused-gap
interpolation, permits edge holding and unbounded gap length, and expresses
several spatial and temporal settings in angular units or milliseconds.
Consequently, a finite prepared coordinate can be temporally filled rather
than directly supported by camera observations. Observed masks, gap lengths,
and fitted-pose support must accompany accuracy and coverage statistics.

An additional articulated-hand stage fits a constrained landmark/capsule
model to the fused sequence. The stable landmark-only fitting path does
not use a MANO surface model or dense depth fitting. Its outputs are saved
alongside the original prepared landmarks, allowing fitted and landmark
poses to be compared. Kinematic regularity of this fitted hand is useful
internal evidence, but it is not independent proof that the recovered hand
motion is correct.

\subsection{Palm frame, interaction point, and gripper signal}

The landmark-derived palm frame uses wrist and metacarpophalangeal (MCP)
landmarks. Its longitudinal axis points from the wrist to the middle MCP;
the normal is the normalised cross product of that axis with the
index-to-pinky MCP vector. The remaining axis completes a right-handed
frame. Degenerate or missing geometry requires a validity decision; an
identity-orientation placeholder is not an observed palm rotation.

The robot target need not be located at the wrist. The current pinch-centre
option uses the midpoint of the thumb and index fingertips for target
position, while orientation can come from the selected palm-pose source.
Thus a run can combine a landmark-derived interaction point with an
articulated-fit orientation. Both choices must be retained in provenance.

For parallel-jaw transfer, the continuous gripper decoder uses a
dimensionless pinch ratio rather than treating fingertip distance as the
robot's metric aperture:
\begin{align}
  r_t &=\frac{\norm{\bx_{\mathrm{thumb},t}-\bx_{\mathrm{index},t}}}
                  {\norm{\bx_{\mathrm{wrist},t}-\bx_{\mathrm{middleMCP},t}}},
       \nonumber\\
  g_t^{\mathrm{raw}}&=
       \left[\frac{r_t-r_{\mathrm{closed}}}
                    {r_{\mathrm{open}}-r_{\mathrm{closed}}}\right]_0^1,
  \qquad
  g_t=g_{t-1}+\alpha_g(g_t^{\mathrm{raw}}-g_{t-1}).
  \label{eq:gripper}
\end{align}
Here zero denotes closed and one fully open. Thresholds and smoothing are
decoder parameters; the first valid sample initialises the filter, and a
missing sample holds the previous command. A binary hysteresis decoder is
also available. These mappings are design choices, rather than anatomical
equivalences between a human hand and a robot gripper. As in other
human-to-robot interfaces, the correspondence must match the target
end-effector \cite{opentv,ace}. A scalar opening signal does not specify a
complete dexterous-hand configuration or contact force.

\section{Object Motion and Relative Representation}
\label{sec:objcentric}

\subsection{Segmentation and rigid modelling}

The implemented scene branch uses prompted SAM~2.1 video segmentation
\cite{sam2}. Prompts associate masks with persistent object identities and
roles. An optional empty-background depth plate supports automatic prompt
proposals from foreground components, but these proposals do not establish
which object is manipulated or which surfaces define the task. Prompt
correction and role assignment are recorded inputs.

Calibrated RGB-D lifting converts masks into semantic point observations.
The object modeller builds a temporally supported canonical core, tracks
it with robust trimmed nearest-neighbour alignment and weighted rigid
fitting, then adds supported shell points and refines the track.
Operator-proximity flags suppress shell growth during possible contact,
reducing the risk of incorporating hand points into the object model.
They indicate geometric proximity, not measured contact or a stable grasp.

The output contains object poses, confidence, support, rotation-information
diagnostics, and validity. Partial visibility and object symmetry may leave
rotation poorly constrained even when translation is stable. These
artifacts are geometric models and motion estimates; they are not yet
textured 3DGS assets and do not automatically alter the hand trajectory.

\subsection{Relative tracks and proposed transfer}

For a saved human frame $H$ and object $O$, the object-relative pose is
\begin{equation}
  \bT^O_{H,t}=(\bT^C_{O,t})^{-1}\bT^C_{H,t}.
  \label{eq:objrel}
\end{equation}
The same construction yields object-relative target and achieved robot
tool poses. The current \texttt{object\_relative.npz} artifact is computed
after object modelling and IK. It permits analysis of relative motion and
interaction error; it is not evidence that general object-relative
retargeting drove the plan. Its saved human frame must also be distinguished
from a pinch-centre robot target when those positions differ.

A possible object-relative target construction for embodiment $m$ is
\begin{equation}
  \bT^{B_m,\mathrm{des}}_{E,t}
    =\bT^{B_m}_{O',t}\,\bT^O_{H,t}\,\bT^H_{E,m},
  \label{eq:transfer}
\end{equation}
where $O'$ is a specified target object trajectory and $\bT^H_{E,m}$ is a
declared human-to-tool correspondence. This equation defines a proposed
transfer option, not the current general-purpose IK input. A constant
correspondence is only an approximation for changing human grasps and
different gripper geometries. Altering object placement may also require
changing approach, support, and release conditions. Object-relative
representations are motivated by their ability to express task geometry
\cite{orion,spot}; representation alone does not make a transferred
interaction feasible.

\subsection{Implemented cube-specific pilot}

The inspected revision includes an experimental object-aware planner for
a cube and Robotiq 2F-85. It requires one explicitly identified manipulated
object, supported canonical geometry with suitable side planes, and
sufficient motion/contact evidence. Missing or ambiguous prerequisites
cause an error rather than a silent hand-only fallback.

This planner separates approach, closure, carry, release, and retreat.
During carry it follows object translation with a fixed pickup offset.
Jaw orientation is aligned with supported cube geometry and held through
the carry, rather than trusting uncertain framewise cube rotation. The
opening command is derived from the modelled pad separation and the
specified cube width. This is a restricted geometric grasp hypothesis.
It intentionally modifies parts of the measured human motion, and neither
its closure command nor its imposed object path demonstrates force closure
or successful manipulation. Hand-only and cube-aware plans must therefore
be evaluated as distinct methods.

\section{Offline Trajectory Retargeting}
\label{sec:retarget}

\subsection{Robot model and target construction}

Each target embodiment specifies an arm model, joint ordering and limits,
base placement, reference posture, gripper, mounting adapter, and
human-to-tool correspondence. The current arm registry contains UR3, UR5,
UR10, UR5e, UR10e, and iiwa14 descriptions. A registry entry is not a
cross-embodiment evaluation result. Panda, Gen3, and xArm7 remain candidate
additions until their models and end-effectors are integrated and tested.

The working Robotiq 2F-85 assembly uses the actual gripper geometry.
Its positional task point is the midpoint of the two modelled contact-pad
centres, and orientation is taken from the designated tool frame.
Because the pad positions vary with gripper opening, the positional task
is not a fixed flange offset. Let $\bq_t\in\mathbb R^{n_m}$ contain the
arm joints and let $g_t$ determine the gripper's drive and mimic joints.
Forward kinematics then provides
$\mathbf p_m(\bq_t,g_t)$ and $\mathbf R_m(\bq_t,g_t)$.
Gripper rate profiling precedes the arm optimisation; gripper joints are
known framewise inputs rather than additional arm decision variables.

The selected human position and orientation are transformed from rig to
workspace coordinates, composed with the configured local tool
correspondence, and then expressed in the robot base frame. These
configuration offsets should not be described as a measured anatomical
or hand--eye calibration. The cube-aware mode replaces the corresponding
targets before the same robot optimisation.

\subsection{Pose error and temporal terms}

For target position $\mathbf p_t^*$ and orientation $\mathbf R_t^*$,
the implemented task residual is
\begin{equation}
  \be_t(\bq_t)=
  \begin{bmatrix}
    \sqrt{w_p}\,[\mathbf p_m(\bq_t,g_t)-\mathbf p_t^*]\\
    \sqrt{w_R}\,\Log_{\mathrm{SO}(3)}
       \!\left((\mathbf R_t^*)^{\!\top}\mathbf R_m(\bq_t,g_t)\right)
  \end{bmatrix}.
  \label{eq:residual}
\end{equation}
The rotation logarithm is represented as a three-vector. This is a
weighted stack of translation and rotation errors, rather than the full
$\SE$ logarithm. The weights set the relative scale of metres and radians;
they are not an estimated measurement covariance.

For $\Delta t=1/f$, temporal operators use
\begin{equation}
  (D_1\bq)_t=\frac{\bq_{t+1}-\bq_t}{\Delta t},\qquad
  (D_2\bq)_t=\frac{\bq_{t+2}-2\bq_{t+1}+\bq_t}{\Delta t^2}.
\end{equation}
The solver estimates $f$ from median positive timestamp intervals and
uses a uniform time step. Recordings with substantial irregular timing
therefore require explicit resampling or a separate timing analysis.
Velocity and acceleration penalties couple adjacent frames, and a posture
penalty favours the selected robot reference. Temporal regularisation and
interaction refinement are also examined in related work such as DexFlow
\cite{dexflow}. Here, the discrete penalties encourage smoothness; they do
not establish continuous-time $C^2$ regularity, bounded jerk, or dynamic
feasibility.

\subsection{Confidence and robust data fitting}

Let $c_{i,t}\in[0,1]$ be the prepared landmark confidence and let
$\mathcal P$ contain wrist, index MCP, middle MCP, and pinky MCP.
The default mean-palm evidence rule is
\begin{equation}
  \bar\omega_t=v_t
     \left(\frac{1}{|\mathcal P|}\sum_{i\in\mathcal P}c_{i,t}\right)^\alpha,
  \label{eq:conf}
\end{equation}
where $v_t$ is pose validity. Under the default episode-normalisation
policy, triangulation quality is divided by the maximum quality in the
episode before clipping. This relative normalisation should not be
interpreted as absolute confidence across recordings. A stricter named
profile uses
\begin{equation}
  \bar\omega_t^{\mathrm{all}}
      =\bar\omega_t\prod_{i\in\mathcal P}o_{i,t},
  \label{eq:confstrict}
\end{equation}
where $o_{i,t}$ indicates direct observation. This all-observed rule is an
alternative, not the inspected default. Under the mean rule a partly
interpolated palm can retain positive pose weight.

Pinch-centre targeting further limits evidence by the lower fingertip
confidence. Invalid target frames receive zero data weight. The solver
clips remaining weights to $[0,1]$ and applies a small positive floor to
nonzero weights; the inspected default floor is $0.05$. Consequently,
arbitrarily weak positive evidence is not attenuated without limit.
Position interpolation and rotation Slerp provide finite placeholders
for invalid target rows, but those rows keep zero data weight.

A Huber loss is applied to the norm of the scaled six-dimensional residual.
Its iteratively reweighted least-squares (IRLS) multiplier satisfies
$\gamma_t^2=\min(1,\delta/\max(\norm{\be_t},\varepsilon))$.
Robust fitting reduces the influence of large residuals; it does not
replace the interpolation and filtering already applied during preparation,
and it cannot detect every plausible but incorrect hand pose.

\subsection{Sparse batch solve and geometric constraints}

The current solver optimises the whole arm trajectory using a custom
sparse Gauss--Newton/IRLS procedure. With current stacked joint vector
$\mathbf q$, pose Jacobians $J_t$, stacked reference
$\mathbf q_{\mathrm{ref}}$, and increment $\Delta\mathbf q$, its QP
subproblem minimises
\begin{equation}
\begin{split}
  \frac12\Bigg[&\sum_t\omega_t\gamma_t^2
       \norm{\be_t+J_t\Delta\bq_t}^{2}
       +\lambda_v\norm{D_1(\mathbf q+\Delta\mathbf q)}^2\\
       &+\lambda_a\norm{D_2(\mathbf q+\Delta\mathbf q)}^2
       +\lambda_r\norm{\mathbf q+\Delta\mathbf q-\mathbf q_{\mathrm{ref}}}^2
       +\mu\norm{\Delta\mathbf q}^2\Bigg].
\end{split}
\label{eq:cost}
\end{equation}
This equation states the implemented local least-squares scaling.
The code's separate logged merit uses the Huber data term plus the full
weighted temporal and posture penalties; it is not numerically identical
to the quadratic model in \eqref{eq:cost}. The fixed damping $\mu$ is a
Tikhonov term, rather than an adaptive Levenberg--Marquardt schedule.

Joint positions and velocities are constrained componentwise:
\begin{align}
  \bq_{\min}&\leq\bq_t\leq\bq_{\max},
       \label{eq:cjoint}\\
  -\dot{\bq}_{\max}&\leq(D_1\bq)_t\leq\dot{\bq}_{\max}.
       \label{eq:cvel}
\end{align}
The geometric constraints have the form
\begin{equation}
  h_j(\bq_t,g_t)\geq0,
  \label{eq:ccoll}
\end{equation}
with local linearisation inside each QP. They cover the workspace floor
plane for represented moving bodies, including the adapter and gripper
contact points, and enabled self-collision pairs. The calibrated plane is
transformed into robot coordinates, including for a tilted base.
These checks do not include arbitrary scene obstacles, general
robot--object contact constraints, torque limits, or hard acceleration
and jerk limits.

Pinocchio supplies kinematics and Jacobians, Pink supplies the configured
self-collision barriers \cite{pink}, and \texttt{qpsolvers} invokes the
sparse OSQP backend. Pink does not itself formulate the entire batch
objective. The trajectory starts from the tiled reference posture, normally
the registered robot home, rather than a sequential frame-to-frame IK
warm start. Each QP also limits increment magnitude. Backtracking checks
the nonlinear merit and represented geometric margins once the current
trajectory is geometrically feasible; an initially infeasible trajectory
may require a feasibility-repair step.

The final trajectory is checked again against exact represented floor
and self-collision margins. The inspected floor tolerance is
$0.5$\,mm; it is a numerical acceptance tolerance for the model, not a
physical clearance guarantee. A small accepted increment can trigger the
solver's convergence flag even if tracking error remains large. Convergence,
constraint satisfaction, target-tracking error, and task success are
therefore separate outcomes.

\subsection{Approach segment and matched baseline}

After retargeting, a quintic approach connects the registered home posture
to the first demonstration configuration. Its duration is increased as
needed to satisfy velocity limits, and the concatenated trajectory is
checked for represented joint, velocity, floor, and enabled collision
constraints. The approach is generated robot motion rather than a measured
part of the human episode. It must be separated from demonstration frames
when aligning video, object motion, or accuracy statistics.

An optional comparison path runs causal framewise IK followed by
Savitzky--Golay smoothing. A meaningful comparison uses the same targets,
robot assembly, gripper sequence, and data weights, and reports resulting
constraint violations as well as tracking error. Smoothing a sequential
solution does not by itself preserve its feasibility.

\section{Proposed 3DGS Reconstruction and Robot Rendering}
\label{sec:gs-method}

\subsection{Scene decomposition and asset acquisition}

The proposed appearance branch uses 3DGS for reusable scene appearance
\cite{gaussiansplatting}, combined with separately represented moving
objects and a robot model. This hybrid design follows the general use of
reconstructed assets in human-to-robot data generation
\cite{real2render2real,robopaint}; it is not presented as a previously absent
combination of components.

The initial target is a static workspace background, a small set of rigid
task objects, and one articulated robot. Human masks exclude the
demonstrator from background fitting. An empty-scene scan or additional
views are needed where the human or task objects otherwise hide the
background. Moving objects are reconstructed in canonical coordinates
and positioned using their estimated tracks. A multiview point cloud
or a set of segmentation masks is an input to this process, not already
a trained Gaussian scene.

For the first implementation, the background is to use 3DGS and the robot
is to use its articulated visual meshes. Task objects may use scanned
meshes or rigid Gaussian assets, with the representation fixed for each
comparison. All assets share the calibrated metric frame. Additional
scans, segmentation corrections, and asset preparation contribute to the
data-collection cost even though they do not require per-robot teleoperation.

\subsection{Reconstruction objective and supported views}

A candidate reconstruction objective for scene parameters $\Theta$ is
\begin{equation}
\begin{split}
  \mathcal L(\Theta)=&\lambda_{\mathrm{rgb}}
       \mathcal L_{\mathrm{rgb}}(M\odot\hat I,M\odot I)
       +\lambda_{\mathrm{ssim}}\mathcal L_{\mathrm{ssim}}(\hat I,I;M)\\
       &+\lambda_z\,
       \frac{\sum_{u}M_z(u)\,|\hat z(u)-z(u)|}
            {\sum_u M_z(u)}.
\end{split}
\label{eq:gs-loss}
\end{equation}
Here $M$ selects valid observations of the modelled component, while
$M_z$ additionally excludes invalid depth and depth-boundary regions.
The depth term is evaluated only when valid pixels exist. A consistent
camera-axis depth convention is required for both measured and rendered
depth. The masked structural term must avoid windows crossing excluded
regions. Loss weights, Gaussian initialisation, optimisation schedule,
and any depth regularisation remain experimental choices to be specified
with the implementation; no performance is assumed from this formulation.

Novel cameras are sampled within a declared operating region supported
by the captured geometry. The method does not assume that a small fixed
rig observes every surface. Disoccluded or poorly supported regions are
marked and included in rejection or coverage statistics. Hallucinated
completion, if later introduced, would require a separately identified
method and evaluation.

\subsection{Time-consistent robot compositing}

At demonstration time $t$, the renderer receives the retargeted arm state,
the profiled gripper command, object poses, camera intrinsics and
extrinsics, and appearance parameters. Robot link transforms are computed
from the same model and joint ordering used for the exported state.
Rendering uses the achieved forward-kinematics pose, rather than drawing
the end-effector at an unreachable human target.

The compositor must resolve visibility between robot meshes, moving
objects, and the Gaussian background in a consistent depth convention.
A foreground overlay without scene-depth tests is insufficient. The
chosen treatment of Gaussian opacity and depth must be specified and
checked at contact boundaries and thin structures. RGB, depth, component
masks, and camera metadata should be generated together when supported.

Object motion initially remains a prescribed estimate from the source
episode. Displaying an object next to the gripper does not establish that
the gripper can cause that motion. Penetration and contact-region checks
can reject geometric inconsistencies, but physical manipulation success
would require a contact simulation or hardware experiment outside this
study's evaluation scope.

\section{Factorial Generation and Episode Acceptance}
\label{sec:factorial-method}

For source episode $d$, let $m\in\{1,\ldots,M\}$ denote embodiment,
$k\in\{1,\ldots,K\}$ a camera configuration, and
$p\in\{1,\ldots,P\}$ an appearance configuration. The proposed output is
a sequence
\begin{equation}
  \mathcal E_d^{m,k,p}=
  \left\{\hat I_{d,t}^{m,k,p},\bq_{d,t}^{m},g_{d,t}^{m},
         \bT^C_{O,d,t},\tau_{d,t},\mathcal V_{d,t}^{m,k,p}\right\}_{t=1}^{T_d},
  \label{eq:factorial-episode}
\end{equation}
where $\mathcal V$ stores relevant validity and support information.
Embodiment-specific retargeting is performed before varying cameras and
appearance, so the same accepted robot trajectory can support multiple
visual observations. Appearance variation is initially restricted to
declared edits that preserve geometry and visibility checks. A lighting
change is not automatically physically valid for a Gaussian model with
baked appearance. Any edit that changes task geometry requires new
retargeting and interaction checks.

The Cartesian product gives $M K P$ requested configurations, not a
guaranteed number of usable demonstrations. Define
$A_d^{m,k,p}\in\{0,1\}$ for configurations that pass the stated checks;
then
\begin{equation}
  N_{d,\mathrm{accepted}}
     =\sum_{m=1}^{M}\sum_{k=1}^{K}\sum_{p=1}^{P}A_d^{m,k,p}
     \leq M K P.
  \label{eq:accepted-expansion}
\end{equation}
Acceptance policies should be fixed before the final comparison and retain
failure reasons. They must distinguish insufficient hand/object evidence,
robot constraints or unacceptable tracking error, unsupported camera
views, visible geometric interaction defects, and inconsistent data
alignment. Thresholds are experimental parameters, not values inferred
from the existence of a processing stage. Current trajectory-export
eligibility is weaker than this proposed synthetic-episode acceptance.

All variants of a source episode share its original motion evidence and
are statistically dependent. Dataset splits must therefore be made by
source episode, and where appropriate by subject, object, or scene,
before factorial expansion. Viewpoint or appearance changes increase
observation variation; they do not create independent manipulation
strategies. Counts, yield, processing time, and manual preparation effort
are reported per source and per factor combination.

\section{Evaluation Protocol}
\label{sec:hocap-method}

\subsection{External hand and object references}

HO-Cap supplies an external multiview RGB-D reference for perception
\cite{hocap}. The protocol specified for the selected 28-sequence
single-hand pick-and-place cohort uses all eight RealSense views and
excludes the HoloLens stream. This cohort is a task/handedness subset,
not an official held-out split or the full dataset. Reproducing its
results requires the exact sequence list, dataset revision, frame counts,
profile manifest, and completed-stage records. The inspected application
source alone does not establish that a benchmark run was completed.

Reference hand labels are excluded from detection, reconstruction, and
parameter fitting. Predictions and references are compared in a common
metric frame by exact frame identity, without root alignment,
Procrustes fitting, temporal search, or reference-based sample selection.
The hand measures are MPJPE, PCK@20/50\,mm, P95 error, and prediction
coverage. Directly observed, interpolated, and articulated-fit outputs
are distinguished. Matched-joint-weighted aggregate means and PCK values
are reported alongside per-sequence statistics; averaging sequence
quantiles does not produce a pooled P95. HO-Cap's semi-automatic
MANO-based annotations are reference labels, not independent motion-capture
ground truth or a competing estimator run inside \viki.

The reference-assisted object protocol supplies published instance masks
and depth to the tracker. It therefore tests tracking conditional on
segmentation, separately from the autonomous SAM-based branch. One
constant object-local transform is fitted at the first confident frame
using the reference pose; subsequent translation and geodesic rotation
errors measure relative tracking drift under that initial alignment.
The documented confidence threshold is $0.35$. Reference-assisted initial
alignment does not measure absolute pose localisation. Coverage,
stationary/moving subsets, and object symmetries must be stated; a raw
geodesic score without symmetry handling has that limitation.

Benchmark error and local-rig self-consistency remain separate evidence.
Neither a small calibration residual nor agreement with a fitted hand
model substitutes for comparison with an independent reference.

\subsection{Retargeting and interaction comparisons}

For each actually tested robot assembly, the evaluation records base
placement, model revision, tool correspondence, solver settings, and
source profile. It measures TCP position and rotation error, accepted
frames and episodes, joint and velocity-limit violations, and represented
floor/self-collision margins. Error is reported separately on
evidence-supported and filled targets. Solver convergence flags are
reported alongside these measures, rather than used as their substitute.

Hand-only and cube-aware retargeting use matched source episodes and
robot assemblies. Object-relative TCP error, contact-region alignment,
closure timing, and geometric penetration characterise their interaction
consistency. Gripper commands are compared against the declared decoder
and actual modelled pad gap. These measures do not become grasp-success
rates without an executed interaction. General object-relative transfer
requires its own implementation and matched comparison before a benefit
can be claimed.

Dense source trajectories preserve more timing information than sparse
keyframes, but filters, gap filling, gripper decoding, reference-posture
penalties, and cube-specific orientation replacement can change the
motion. Where natural-motion preservation is investigated, it is measured
through timing, path, and orientation differences on supported intervals;
it is not assumed from markerless capture.

\subsection{Rendering and expansion experiments}

The proposed reconstruction evaluation holds out cameras and scene states
before tuning. PSNR, SSIM, and valid-pixel depth RMSE compare corresponding
scene components in matching states, with background and task-object
regions reported separately. A static background model is not evaluated
against unmasked moving humans. A human demonstration image is also not
a paired ground-truth image of the replacement robot.

Compositing comparisons use identical trajectories, object states, and
camera poses. They inspect depth ordering, robot/object boundaries,
temporal stability, and agreement between rendered geometry and exported
state. Paired perceptual metrics require a corresponding target image;
otherwise the evaluation uses explicitly defined geometric and consistency
tests rather than a misleading full-image similarity score.

Factorial expansion is compared with single-configuration generation on
the same source episodes. The report includes requested and accepted
combinations, rejection categories, supported viewpoint range, runtime,
and manual intervention. Any comparison with teleoperation effort must
include capture, calibration, scans, corrections, robot-model preparation,
computation, and rejected outputs. Policy-training benefit is not inferred
from appearance diversity or episode count.

\section{Export, Provenance, and Replay Status}
\label{sec:export}

The working default exporter writes a self-contained NumPy trajectory
bundle and JSON metadata. It retains joint trajectories and derivatives,
normalised gripper opening, modelled metric opening, target and achieved
poses, tracking errors, source timing, and evidence information. Optional
object-model and post-IK object-relative artifacts are copied when their
stages are complete. Units, frame names, source episode, processing
profile, robot assembly, and parameter provenance must remain available.
In the inspected plan schema, the metric gripper opening is based on the
URDF pad-centre gap; older schemas may use nominal width and must be
interpreted by version.

Current export requires a readable plan, not a demonstrated manipulation
outcome. It may include unlabelled episodes and reports missing or skipped
artifacts in its manifest. The optional LeRobot path has stricter
eligibility and partial annotation support; shape-correct placeholders
are not measured object-relative labels. The existing bundle is therefore
a trajectory artifact, rather than the completed rendered dataset in
\eqref{eq:factorial-episode}.

For the proposed synthetic export, every rendered frame must be paired
with the state that produced it. The schema must distinguish state at
$t$ from a command for $t+1$ and specify whether an action is an absolute
joint target, an increment, or a velocity command. Camera calibration,
appearance seed, renderer and asset versions, validity masks, and
acceptance decisions are additional provenance. Original human-camera
videos cannot be relabelled as robot observations merely because a robot
trajectory has been generated.

\subsection{Physical replay}
\label{sec:replay}

The current dry-run replay path echoes a plan and does not measure robot
tracking. The physical driver and several screening checks remain
unimplemented. No physical-replay acceptance or achieved hardware state
can therefore be inferred from a dry-run result. Hardware execution,
contact simulation, and policy training lie outside the present
reconstruction, retargeting, and rendering evaluation. They would supply
additional evidence needed for claims about executable manipulation or
learning effectiveness.

\section{Assumptions and Operating Limits}
\label{sec:assumptions}

The initial scope is a fixed, calibrated RGB-D rig, one tracked hand,
rigid task objects, a bounded workspace, and parallel-jaw transfer.
Handedness, object identities, task roles, and robot placement must be
declared. Unknown geometry and uncertain observations remain explicit
rather than being treated as valid because an intermediate file exists.

The reconstruction assumes sufficient visibility and stable calibration.
The retargeter assumes that the selected human-to-tool correspondence and
robot model are appropriate for the task. Its collision checks cover only
the represented geometry. The proposed appearance branch assumes a
separable static background and rigid moving assets; deformable objects,
general scene relighting, and arbitrary unseen surfaces require further
models and validation.

The intended outcome is a measured set of geometrically consistent
candidate robot episodes derived without per-robot teleoperation during
human capture. Whether a source episode supports multiple embodiments,
views, and scene appearances is an experimental result. Preservation of
task-relevant human motion, physical manipulability, and downstream
learning benefit are separate claims with separate evidence requirements.